% Source: the_product_structure_in_ML.tex
\begin{notechapter}
\chapter{The Product Structure in Machine Learning}\label{ch:product-structure}
\chaptersubtitle{Likelihood factorization and naive Bayes}

\begin{abstract}
    This document explores the product structures that appear in machine learning, particularly focusing on the relationship between the likelihood function in maximum likelihood estimation and the factorization in Naive Bayes models. While both involve products of probabilities, they represent fundamentally different concepts and assumptions. We clarify these differences and show how they interact when training probabilistic models.
\end{abstract}

\section{Introduction: Two Product Structures}

In machine learning, we frequently encounter products of probabilities. Two common instances are:

\begin{itemize}
    \item The \textbf{likelihood function} in Maximum Likelihood Estimation (MLE): $\mathcal{L}(\theta) = \prod_{i=1}^n Q_\theta(y_i|x_i)$
    \item The \textbf{Naive Bayes classifier}: $P(Y|X) \propto P(Y) \prod_{j=1}^d P(X_j|Y)$
\end{itemize}

At first glance, both involve products, but they serve different purposes and make different assumptions. This document clarifies their relationship and distinctions.

\section{Maximum Likelihood Estimation and the Likelihood Product}

\subsection{The Likelihood Function}

Given a dataset $D = \{(x_1, y_1), (x_2, y_2), \ldots, (x_n, y_n)\}$ and a model with parameters $\theta$ that outputs a distribution $Q_\theta(Y|X)$, the likelihood is:

\begin{definition}[Likelihood]
\[
\mathcal{L}(\theta) = \prod_{i=1}^n Q_\theta(y_i | x_i)
\]
\end{definition}

The log-likelihood is:
\[
\log \mathcal{L}(\theta) = \sum_{i=1}^n \log Q_\theta(y_i | x_i)
\]

\subsection{Assumptions and Interpretation}

The product structure $\prod_{i=1}^n$ comes from the \textbf{i.i.d. (independent and identically distributed)} assumption:

\begin{itemize}
    \item \textbf{Independence}: Data points $(x_i, y_i)$ are independent of each other
    \item \textbf{Identical distribution}: All data points come from the same underlying distribution
\end{itemize}

This is a \textbf{reasonable and standard assumption} in most supervised learning scenarios.

\section{Naive Bayes and Its Product Structure}

\subsection{The Naive Bayes Model}

Naive Bayes is a probabilistic classifier based on Bayes' theorem with a strong independence assumption:

\begin{definition}[Naive Bayes Classifier]
For a data point with features $X = (X_1, X_2, \ldots, X_d)$ and class $Y$:
\[
P(Y|X) = \frac{P(Y) \prod_{j=1}^d P(X_j|Y)}{P(X)} \propto P(Y) \prod_{j=1}^d P(X_j|Y)
\]
\end{definition}

\subsection{Assumptions and Interpretation}

The product structure $\prod_{j=1}^d$ comes from the \textbf{conditional independence assumption}:

\begin{itemize}
    \item \textbf{Conditional independence}: Features $X_1, X_2, \ldots, X_d$ are conditionally independent given the class $Y$
    \item \textbf{Mathematically}: $P(X_1, X_2, \ldots, X_d | Y) = \prod_{j=1}^d P(X_j | Y)$
\end{itemize}

This is a \textbf{strong and often incorrect assumption}, but it works surprisingly well in practice.

\section{Key Differences Between the Two Products}

\begin{table}[h]
\centering
\begin{tabular}{p{0.2\textwidth}p{0.35\textwidth}p{0.35\textwidth}}
\toprule
\textbf{Aspect} & \textbf{MLE Likelihood Product} & \textbf{Naive Bayes Product} \\
\midrule
\textbf{What multiplies} & Data points across dataset & Features within a data point \\
\hline
\textbf{Notation} & $\prod_{i=1}^n$ over $n$ instances & $\prod_{j=1}^d$ over $d$ features \\
\hline
\textbf{Assumption} & Data point independence (i.i.d.) & Feature conditional independence \\
\hline
\textbf{Purpose} & Training principle & Model structure \\
\hline
\textbf{Scope} & General methodology & Specific model class \\
\hline
\textbf{Reasonableness} & Standard, reasonable assumption & Strong, often violated assumption \\
\bottomrule
\end{tabular}
\caption{Comparison of the two product structures}
\end{table}

\section{When They Meet: Training Naive Bayes with MLE}

The connection becomes clear when we \textbf{train} a Naive Bayes model using maximum likelihood estimation.

\subsection{Complete Likelihood for Naive Bayes}

When training Naive Bayes, we combine both product structures:

\[
\mathcal{L}(\theta) = \prod_{i=1}^n \left[ P(y_i) \prod_{j=1}^d P(x_{i,j}|y_i) \right]
\]

This combines:
\begin{itemize}
    \item Product over data points: $\prod_{i=1}^n$ (from MLE/i.i.d. assumption)
    \item Product over features: $\prod_{j=1}^d$ (from Naive Bayes conditional independence)
\end{itemize}

\subsection{Training Procedure}

The log-likelihood becomes:
\[
\log \mathcal{L}(\theta) = \sum_{i=1}^n \left[ \log P(y_i) + \sum_{j=1}^d \log P(x_{i,j}|y_i) \right]
\]

We can maximize this by estimating:
\begin{itemize}
    \item Class priors: $\hat{P}(Y=y) = \frac{\text{count}(Y=y)}{n}$
    \item Feature distributions: $\hat{P}(X_j=x|Y=y) = \frac{\text{count}(X_j=x \text{ and } Y=y)}{\text{count}(Y=y)}$
\end{itemize}

\section{Concrete Example: Email Spam Classification}

Let's consider classifying emails as spam/ham based on presence of words "money" and "free".

\subsection{Naive Bayes Prediction}

For a single email containing both "money" and "free":
\[
P(\text{spam}|\text{money, free}) \propto P(\text{spam}) \times P(\text{money}|\text{spam}) \times P(\text{free}|\text{spam})
\]

\textbf{Product over features} for one data point.

\subsection{MLE Training}

For our entire dataset of $n$ emails:
\[
\mathcal{L} = \prod_{i=1}^n \left[ P(y_i) \times P(\text{money}_i|y_i) \times P(\text{free}_i|y_i) \right]
\]

\textbf{Product over data points} across the dataset.

\begin{remark}[Joint versus conditional likelihood]
The likelihood maximized when training Naive Bayes is the \emph{joint} likelihood \( \prod_i P(x_i, y_i) = \prod_i P(y_i)\prod_j P(x_{i,j}\mid y_i) \): Naive Bayes is a \emph{generative} model. This is not the conditional likelihood \( \prod_i Q_\theta(y_i\mid x_i) \) introduced at the beginning of this chapter, which is what discriminative models (logistic regression, neural networks) maximize. The closed-form count estimates above are the maximizers of the joint likelihood; maximizing the conditional likelihood of the same model has no closed form.
\end{remark}

\section{General Perspective}

\subsection{MLE with Cross-Entropy as Universal Principle}

Maximum likelihood estimation is a \textbf{general training principle} applicable to:

\begin{itemize}
        \item Naive Bayes models (joint likelihood \( \prod_i P(x_i, y_i) \))
    \item Logistic regression (conditional likelihood \( \prod_i Q_\theta(y_i \mid x_i) \), i.e.\ cross-entropy)
    \item Neural networks with softmax output
    \item Any probabilistic classifier
\end{itemize}

\subsection{Naive Bayes as Specific Model}

Naive Bayes is a \textbf{specific model class} characterized by:
\begin{itemize}
    \item Conditional independence assumption
    \item Simple parameter estimation
    \item Fast training and inference
    \item Often used as a baseline model
\end{itemize}

\section{Mathematical Relationship}

\subsection{From Naive Bayes to Logistic Regression}

Interestingly, if we assume Gaussian distributions for features in Naive Bayes, \emph{with a variance \( \sigma_j^2 \) shared by the two classes}, the posterior has the logistic regression form:

For binary classification with Naive Bayes and Gaussian features:
\[
\log\frac{P(Y=1|X)}{P(Y=0|X)} = \log\frac{P(Y=1)}{P(Y=0)} + \sum_{j=1}^d \log\frac{\mathcal{N}(X_j;\mu_{j1},\sigma_j^2)}{\mathcal{N}(X_j;\mu_{j0},\sigma_j^2)}
\]

This simplifies to:
\[
\log\frac{P(Y=1|X)}{P(Y=0|X)} = w_0 + \sum_{j=1}^d w_j X_j
\]

Which is exactly the logistic regression form, with \( w_j = (\mu_{j1} - \mu_{j0})/\sigma_j^2 \) and \( w_0 = \log\frac{P(Y=1)}{P(Y=0)} + \sum_j \frac{\mu_{j0}^2 - \mu_{j1}^2}{2\sigma_j^2} \). The two models have the same form of \( P(Y \mid X) \) but are \emph{trained differently}: Gaussian Naive Bayes estimates the means and variances by maximizing the joint likelihood, whereas logistic regression fits \( w \) directly by maximizing the conditional likelihood. They form a generative--discriminative pair: logistic regression is asymptotically better when the conditional independence assumption is wrong, while Naive Bayes often converges faster with little data (Ng and Jordan, 2002). With class-dependent variances the log-odds become quadratic in \( X \).

\subsection{The Complete Picture}

\begin{figure}[h]
\centering
\begin{tabular}{c}
\textbf{Training Principle} \\
Maximum Likelihood Estimation \\
$\mathcal{L}(\theta) = \prod_{i=1}^n Q_\theta(y_i|x_i)$ \\
$\Downarrow$ \\
\textbf{Model Classes} \\
\begin{tabular}{ccc}
Naive Bayes & Logistic Regression & Neural Networks \\
$P(Y)\prod_j P(X_j|Y)$ & $\sigma(w^T X)$ & $\text{softmax}(f_\theta(X))$ \\
\end{tabular} \\
$\Downarrow$ \\
\textbf{Loss Function} \\
Negative log-likelihood: joint $-\frac{1}{n}\sum_{i} \log P_\theta(x_i, y_i)$ (Naive Bayes)\\ or conditional (cross-entropy) $-\frac{1}{n}\sum_{i} \log Q_\theta(y_i|x_i)$ (discriminative models)
\end{tabular}
\caption{Hierarchy from training principles to specific models}
\end{figure}

\section{Empirical Considerations}

\subsection{When to Use Each Approach}

\textbf{Naive Bayes is preferred when:}
\begin{itemize}
    \item Dataset is small
    \item Training time is critical
    \item Features are approximately conditionally independent
    \item You need a simple, interpretable baseline
\end{itemize}

\textbf{MLE with flexible models (neural networks) is preferred when:}
\begin{itemize}
    \item Large datasets are available
    \item Computational resources are sufficient
    \item Complex feature interactions exist
    \item Maximum predictive performance is desired
\end{itemize}

\subsection{Practical Implementation}

In practice, both approaches minimize a negative log-likelihood, but not the same one: Naive Bayes minimizes \( -\sum_i \log P(x_i, y_i) \) (in closed form, by counting), while discriminative models minimize the cross-entropy \( -\sum_i \log Q(y_i \mid x_i) \) (numerically). The sketch below shows the discriminative version:

\begin{example}[Python Pseudocode]
\begin{verbatim}
# Both approaches minimize negative log-likelihood
def negative_log_likelihood(params, data):
    total_loss = 0
    for x, y_true in data:
        
        # Neural Net: compute Q(Y|X) using forward pass
        prob = model_prediction(x, params)
        total_loss += -log(prob[y_true])
    return total_loss

# Then optimize parameters
optimal_params = minimize(negative_log_likelihood, initial_params)
\end{verbatim}
\end{example}

\section{Conclusion}

The product structure $\prod_{i=1}^n Q_\theta(y_i|x_i)$ in maximum likelihood estimation and the product structure $\prod_{j=1}^d P(X_j|Y)$ in Naive Bayes represent different but complementary concepts:

\begin{itemize}
    \item The \textbf{MLE product} represents the assumption that data points are independent, which is fundamental to most statistical learning approaches.
    
    \item The \textbf{Naive Bayes product} represents a modeling assumption about feature independence within each data point, which is specific to the Naive Bayes model class.
    
    \item When we train Naive Bayes models, we combine both product structures: we assume data points are independent \textbf{and} features are conditionally independent.
    
        \item Maximum likelihood estimation provides a unified framework that encompasses both simple generative models like Naive Bayes (joint likelihood) and discriminative models like logistic regression and neural networks (conditional likelihood, i.e.\ cross-entropy).
\end{itemize}

Understanding this distinction helps clarify the theoretical foundations of machine learning and explains why cross-entropy loss is so ubiquitous across different model architectures.
\end{notechapter}
